% Reviewer 2
\reviewer
\begin{revcomment}[2.1]
	Thanks for the authors’ efforts in revising the paper. All my comments are well addressed. There only remain minor editorial issues that shall be tackled when finalizing the paper:
\begin{itemize}
    \item Abstract: The last sentence, ``Supply-demand dynamics also affect the system performance'', is too vague and should summarize the findings more specifically.
    \item Page 4: The second paragraph of contributions is dedicated to the greedy mechanism. Hence, the last sentence, ``Furthermore, the developed two mechanisms \dots'', may be better moved to the next paragraph or simply removed.
    \item Page 24: The notation $\pi$ in Proposition 9 and in the following discussion is not clearly defined.
    \item Page 35: The explanation. ``while the utility loss is bound to happen\dots'', is not very clear (as least to me).
\end{itemize}
\end{revcomment}

\hypertarget{Response to R2}{}
\begin{revresponse}[2.1]
	We sincerely thank Reviewer 2 for the positive evaluation and for confirming that the previous comments have been securely addressed. We have carefully incorporated all remaining editorial suggestions within this final revision. 
    \begin{itemize}
        \item Abstract: We have revised the abstract to explicitly detail how variations in supply-demand ratios directly invert the comparative profitability margins and fulfillment characteristics between the VCG and greedy mechanisms. 
        \item Page 4: We agree with this structural refinement and we have moved the summarizing comparative sentence to the subsequent paragraph to maintain narrative flow logically. 
        \item Page 24: We have now formally defined the parameter $\pi$ within the context of Proposition 9, explicitly outlining its mathematical representations and physical significance. 
        \item Page 35: We have completely rewritten this explanation to state the root cause transparently: rather than arbitrary losses, the utility loss inherent to approximation mechanisms acts as a deterministic mathematical outcome of artificially bounding the combinatorial search tree depth. We detail the mechanism showing exactly how early iterations trigger this trade-off.
    \end{itemize}
\end{revresponse}



